\subsection{多项式凸集}

定义 3. --- 设 \(\Lambda\) 是一个集合，V 是 \(\mathbf{C}^{\Lambda}\) 的一个子集。我们说 V 是多项式凸的，如果 V 是点 \((c_{\lambda})_{\lambda \in \Lambda}\)（\(\mathbf{C}^{\Lambda}\) 中的点）所成的集合，这些点满足

\[
| \mathrm{P} ((c _ {\lambda})) | \leqslant \sup _ {c \in \mathrm{V}} | \mathrm{P} (c) |
\]

对一切 \(\mathrm{P}\in \mathbf{C}[(\mathrm{X}_{\lambda})]\) 成立。

引理 3. --- 设 \(\Lambda\) 是一个集合。\(\mathbf{C}^{\Lambda}\) 的一个子集 V 是多项式凸的，当且仅当存在一个族 \((\mathrm{P}_i)_{i\in \mathrm{I}}\)（由 \(\mathbf{C}[(\mathrm{X}_{\lambda})_{\lambda \in \Lambda}]\) 的元素组成）与一个族 \((\mathrm{M}_i)_{i\in \mathrm{I}}\)（由 \([0, +\infty]\) 的元素组成），使得 V 是 \(c\in \mathbf{C}^{\Lambda}\) 中满足

\[
| \mathrm{P} _ {i} (c) | \leqslant \mathrm{M} _ {i}
\]

（对一切 i ∈ I）的元素所成的集合。

若 \(C^{\Lambda}\) 的子集 V 是多项式凸的，则它对由 \(\mathbf{C}[(X_{\lambda})_{\lambda\in\Lambda}]\) 的元素 P 组成的族满足上述条件，其中取 \(M_{P}=\sup_{c\in V}|P(c)|\)。

反之，设 \((\mathrm{P}_{i})_{i\in\mathrm{I}}\) 是由 \(\mathbf{C}[(\mathrm{X}_{\lambda})_{\lambda\in\Lambda}]\) 的元素组成的一个族，并设 \((\mathrm{M}_{i})_{i\in\mathrm{I}}\) 是由 \([0,+\infty]\) 的元素组成的一个族。设 V 是 \(C^{\Lambda}\) 中使得 \(|\mathrm{P}_{i}(c)|\leqslant\mathrm{M}_{i}\)（对一切 \(i\in I\)）的 c 所成的集合。于是我们有 \(\sup_{c\in\mathrm{V}}|\mathrm{P}_{i}(c)|\leqslant\mathrm{M}_{i}\) 对 \(i\in I\) 成立。假定 \(x\in C^{\Lambda}\) 满足

\[
| \mathrm{P} (x) | \leqslant \sup _ {c \in \mathrm{V}} | \mathrm{P} (c) |
\]

对一切 \(\mathrm{P}\in\mathbf{C}[(\mathrm{X}_{\lambda})]\) 成立。特别地，对 \(i\in I\)，我们有 \(|\mathrm{P}_{i}(x)|\leqslant\mathrm{M}_{i}\)，因而 \(x\in V\)。反过来，对每个元素 \(x\in V\) 及每个多项式 \(\mathrm{P}\in\mathbf{C}[(\mathrm{X}_{\lambda})]\)，有 \(|\mathrm{P}(x)|\leqslant\sup_{c\in\mathrm{V}}|\mathrm{P}(c)|\)。因此，集合 V 是多项式凸的。

引理 4. --- 设 A 是一个交换幺 Banach 代数。设 \(\Lambda\) 是一个集合，\(x = (x_{\lambda})_{\lambda \in \Lambda}\) 是 A 的元素的一个族。若由族 x 生成的幺子代数在 A 中稠密，则联合谱 \(\mathrm{Sp}_{\mathrm{A}}^{\Lambda}(x)\) 是多项式凸的。

这由 I, p.~43, 命题 12 的断言 \(b\) ) 以及定义 3 得出。

\(C^{\Lambda}\) 的多项式凸子集的任意交是多项式凸的（引理 3）。这证明了下述定义的合理性：

定义 4. --- 设 \(\Lambda\) 是一个集合，V 是 \(\mathbf{C}^{\Lambda}\) 的一个子集。V 的多项式凸包是 \(\mathbf{C}^{\Lambda}\) 中包含 V 的最小多项式凸子集。

V 的多项式凸包是 \(C^{\Lambda}\) 中满足 \(|P(c)| \leqslant \sup_{W} |P|\)（对一切 \(P \in \mathbf{C}[(X_{\lambda})]\)）的 c 所成的集合。事实上，依引理 3 这个集合是多项式凸的，且依定义它含于每个包含 V 的多项式凸集中。

例. --- 设 \(\Lambda\) 是一个有限集合，\(V \subset \mathbf{C}^{\Lambda}\) 是一个紧凸子集。则 \(V\) 是多项式凸的。事实上，设 \(W\) 是 \(V\) 的多项式凸包。我们来证明 \(W \subset V\)，由此即得该断言。设 \(x \in \mathbf{C}^{\Lambda} - V\)。存在一个实超平面 \(H\)（\(\mathbf{C}^{\Lambda}\) 中的实超平面），它严格分离 \(x\) 与 \(V\) (EVT, II, p.~41, 命题 4)。设 \(f_{\mathbf{R}}\) 是一个 \(\mathbf{R}\) -线性形式，定义在 \(\mathbf{C}^{\Lambda}\) 上，并设 \(\alpha \in \mathbf{R}\)，使得 \(H\) 是 \(y \in \mathbf{C}^{\Lambda}\) 中满足 \(f_{\mathbf{R}}(y) = \alpha\) 的 y 所成的集合。设 \(f\) 是 \(\mathbf{C}^{\Lambda}\) 上的一个线性形式，使得 \(f_{\mathbf{R}} = \mathcal{R}(f)\)。于是我们有

\[
\mathcal {R} (f (x)) > \sup _ {y \in V} \mathcal {R} (f (y)).
\]

对每个 \(t \in R\) 与 \(y \in V\)，令 \(f_{t}(y) = t + f(y)\)。我们有 \(|f_{t}| - \mathcal{R}(f_{t}) \to 0\)（在赋以紧收敛拓扑的 \(\mathcal{C}(\mathbf{C}^{\Lambda}, \mathbf{R})\) 中，当 \(t \to +\infty\) 时）。对充分大的 t，由于 V 是紧的，由此得出 \(|f_{t}(x)| > \sup_{y \in V} |f_{t}(y)|\)。于是 \(x \in C^{\Lambda} - W\)，因为 \(f_{t}\) 是多项式函数。

引理 5. --- 设 K 是 C 的一个紧子集。我们用 \(\widehat{\mathbf{K}}\) 表示 K 与 \(\mathbf{C} - \mathbf{K}\) 的相对紧连通分支的并。则集合 \(\widehat{\mathbf{K}}\) 是紧的。

由于 K 是紧的，存在实数 r > 0，使得 K 含于以 0 为中心、以 r 为半径的开圆盘 D。于是 C-K 包含 C-D。由于空间 C-D 是连通的（它同胚于 \([r, +\infty[\times\mathbf{S}^1)\)），它含于 \(\mathbf{C}-\mathbf{K}\) 的一个连通分支 U。\(\mathbf{C}-\mathbf{K}\) 的任何其他连通分支都含于 D，因而都是有界的。于是连通分支 U 是 \(\mathbf{C}-\mathbf{K}\) 的唯一的无界连通分支，即 \(\mathrm{U} = \mathbf{C}-\widehat{\mathrm{K}}\)。由于 U 是开的且包含圆盘 D 的余集，集合 \(\widehat{\mathrm{K}}\) 是紧的。

命题 13. --- 设 \(n \geqslant 1\) 是一个整数。设 \(\mathbf{K} \subset \mathbf{C}^{n}\) 是一个闭子集，V 是它的多项式凸包。

\begin{enumerate}
\def\labelenumi{\alph{enumi})}
\item
  \(C^{n} - K\) 的每个有界连通分支都含于 V；
\item
  若 n = 1 且 K 是紧的，则 V 是 K 与 C-K 的有界连通分支的并 \(\widehat{K}\)。
\end{enumerate}

由于 \(K \subset C^{n}\) 是闭的，它的余集 \(C^{n}-K\) 是开的，从而是局部连通的，因此 \(C^{n}-K\) 的每个连通分支都是开的。最大模原理 (VAR, R1, p.~29, 3.3.7) 表明 \(C^{n}-K\) 的每个有界连通分支都含于 V，这就证明了断言 a)。

现在假定 \(n = 1\) 且 K 是紧的。用 \(\widehat{\mathbf{K}}\) 表示 K 与 \(\mathbf{C} - \mathbf{K}\) 的有界连通分支的并，于是 \(\widehat{\mathbf{K}} \subset \mathbf{V}\)；由前所述，集合 \(\widehat{\mathbf{K}}\) 是紧的（引理 5）。设 A 是交换幺 Banach 代数 \(\mathcal{C}(\widehat{\mathbf{K}})\)。设 \(x \in \mathbf{A}\) 是 \(\widehat{\mathbf{K}}\) 上的恒等函数，B 是 A 的由 \(x\) 生成的闭幺子代数。我们有 \(\mathrm{Sp}_{\mathrm{A}}(x) = \widehat{\mathbf{K}}\) (I, p.~17, 例 3)，因而 \(\mathbf{C} - \mathrm{Sp}_{\mathrm{A}}(x)\) 是连通的。于是我们有 \(\mathrm{Sp}_{\mathrm{B}}(x) = \mathrm{Sp}_{\mathrm{A}}(x)\) (I, p.~28, 命题 6 的推论)。由于依引理 4 \(\mathrm{Sp}_{\mathrm{B}}(x)\) 是多项式凸的，这表明 \(\widehat{\mathbf{K}}\) 是多项式凸的，从而 \(\mathbf{V} \subset \widehat{\mathbf{K}}\)。

该命题的第二部分不能推广到 \(n \geqslant 2\) 的情形（参见 I, p.~170, 习题 23）。

命题 14. --- 设 \(\Lambda\) 是一个集合。设 A 是一个交换幺 Banach 代数，\(x = (x_{\lambda})_{\lambda \in \Lambda}\) 是 A 的元素的一个族，A' 是由 \(x\) 生成的幺 Banach 子代数。则联合谱 \(\mathrm{Sp}_{\mathrm{A}'}^{\Lambda}(x)\) 是 \(\mathrm{Sp}_{\mathrm{A}}^{\Lambda}(x)\) 的多项式凸包。

事实上，I, p.~43, 命题 12 b) 证明了 \(\mathrm{Sp}_{\mathrm{A}'}^{\Lambda}(x)\) 是由 \(c \in \mathbf{C}^{\Lambda}\) 中满足 \(|\mathrm{P}(c)| \leqslant \varrho(\mathrm{P}(x))\)（对一切 \(\mathrm{P} \in \mathbf{C}[(\mathrm{X}_{\lambda})]\)）的 c 所成的集合。

现在我们有

\[
\varrho (\mathrm{P} (x)) = \sup _ {\chi \in \mathsf {X} (\mathrm{A})} | \chi (\mathrm{P} (x)) | = \sup _ {\chi \in \mathsf {X} (\mathrm{A})} | \mathrm{P} ((\chi (x _ {\lambda})) _ {\lambda \in \Lambda}) | = \sup _ {c \in \operatorname{Sp} _ {\mathrm{A}} ^ {\Lambda} (x)} | \mathrm{P} (c) |,
\]

这是依 I, p.~38, 命题 7 a) 及 I, p.~43, 命题 12 a)。于是由引理 3 即得结论。

命题 15. --- 设 \(\Lambda\) 是一个集合。设 K 是 \(\mathbf{C}^{\Lambda}\) 的一个紧多项式凸子集。设 \(\mathrm{A}_0\) 是 \(\mathbf{C}^{\Lambda}\) 上的多项式函数在 K 上的限制所成的集合，并设 A 是 \(\mathrm{A}_0\) 在代数 \(\mathcal{C}(\mathrm{K})\) 中的闭包。

设 ev: K → X(A) 是由 x ↦ ev\_x 定义的映射，其中 ev\_x 是 A 的特征标 f ↦ f(x)。设 z = (z\_λ)\_\{λ∈Λ\} 是 A 中由 C\^{}Λ 上的坐标函数在 K 上的限制所成的族，并设 φ 是由族 z 定义的、从 X(A) 到 Sp\_A\^{}Λ(z) 上的满射。

于是我们有 K = Sp \(_{A}^{\Lambda}(z)\)，且映射 ev: K → X(A) 与 φ: X(A) → K 是互逆的同胚。

映射 \(\varphi \circ\) ev 是 K 的恒等映射。特别地，K 含于像 \(\mathrm{Sp}_{\mathrm{A}}^{\Lambda}(z)\)（\(\varphi\) 的像）之中。

命题 14 蕴涵 \(\mathrm{Sp}_{\mathrm{A}}^{\Lambda}(z)\) 是 \(\mathrm{Sp}_{\mathcal{C}(\mathrm{K})}^{\Lambda}(z)\) 的多项式凸包。另一方面，由于我们有 \(\mathrm{Sp}_{\mathcal{C}(\mathrm{K})}^{\Lambda}(z) = \mathrm{K}\) (I, p.~42, 例)，且依假设它是多项式凸的，我们推出 \(\mathrm{Sp}_{\mathrm{A}}^{\Lambda}(z) = \mathrm{K}\)。

由于由元素族 \(z_{\lambda}\) 生成的幺子代数在 A 中稠密，映射 \(\varphi\) 是从 \(\mathsf{X}(\mathsf{A})\) 到 \(\mathrm{Sp}_{\mathsf{A}}^{\Lambda}(z)=\mathrm{K}\) 上的一个同胚（第 I 卷, p. 43, 命题 12(a)）。恒等式 \(\varphi\circ ev=Id_{K}\) 于是表明 ev 是 \(\varphi\) 的逆同胚。

对每个子集 \(\Lambda'\)（\(\Lambda\) 的子集），我们用 \(\mathrm{pr}_{\Lambda'}\) 表示典范投影 \(\mathbf{C}^{\Lambda} \to \mathbf{C}^{\Lambda'}\)。设 W 是 \(\mathbf{C}^{\Lambda}\) 的一个子集，V 是它的多项式凸包。令 \(W' = \mathrm{pr}_{\Lambda'} W\)。由于 \(\mathbf{C}[(X_{\lambda})_{\lambda \in \Lambda'}]\) 的每个元素都与 \(\mathbf{C}[(X_{\lambda})_{\lambda \in \Lambda}]\) 的一个元素视为同一，\(W'\) 的多项式凸包含于 \(\mathrm{pr}_{\Lambda'} V\)。

引理 6. --- 设 \(K \subset C^{\Lambda}\) 是一个紧多项式凸子集，U 是 K 的一个邻域。存在一个有限子集 \(\Lambda_{0}\)（\(\Lambda\) 的子集），使得对每个子集 \(\Lambda'\)（\(\Lambda\) 的包含 \(\Lambda_{0}\) 的子集），集合 \(\mathrm{pr}_{\Lambda'}(\mathrm{U})\) 包含 \(\mathrm{pr}_{\Lambda'}(\mathrm{K})\) 的多项式凸包。

由于 K 是紧的，存在 C 中一族紧圆盘 \(D_{\lambda}\)（以 0 为中心，半径为 \(R_{\lambda}\)），使得 K 含于乘积 \(\mathrm{D} = \prod_{\lambda} \mathrm{D}_{\lambda}\)。对每个 \(\mathrm{P} \in \mathbf{C}[(\mathrm{X}_{\lambda})]\)，用 \(\mathrm{K}_{\mathrm{P}}\) 表示 \(x \in \mathbf{C}^{\Lambda}\) 中满足

\[
| \mathrm{P} (x) | \leqslant \sup _ {c \in \mathrm{K}} | \mathrm{P} (c) |.
\]

我们有

\[
\mathrm{D} \cap \bigcap_ {\mathrm{P}} \mathrm{K} _ {\mathrm{P}} = \mathrm{K}\tag{3}
\]

它含于 U。于是，空间 D 是开集 \(D \cap U\) 与诸开集 \(D - K_{P}\)（\(P \in \mathbf{C}[(X_{\lambda})]\)）的并。由于 D 是紧的，存在一个整数 \(q \geqslant 1\) 及多项式 \(P_{1}, \ldots, P_{q} \in \mathbf{C}[(X_{\lambda})]\)，使得：

\[
\mathrm{D} \cap \mathrm{K} _ {\mathrm{P} _ {1}} \cap \dots \cap \mathrm{K} _ {\mathrm{P} _ {q}} \subset \mathrm{U}.\tag{4}
\]

存在一个有限集 \(\Lambda_0 \subset \Lambda\)，使得 \(P_i \in \mathbf{C}[(X_\lambda)_{\lambda \in \Lambda_0}]\) 对 \(1 \leqslant i \leqslant q\) 成立。我们来证明 \(\Lambda_0\) 具有所要求的性质。

设 \(\Lambda'\) 是 \(\Lambda\) 的一个包含 \(\Lambda_0\) 的子集。设 E 是 \(\mathbf{C}^{\Lambda'}\) 的子集，由元素 \(c = (c_{\lambda})_{\lambda \in \Lambda'}\)（由不等式 \(|c_{\lambda}| \leqslant \mathrm{R}_{\lambda}\)（对 \(\lambda \in \Lambda'\)）所定义）组成，且

\[
| \mathrm{P} _ {i} (c) | \leqslant \sup _ {c \in \mathrm{K}} | \mathrm{P} _ {i} (c) |
\]

（对 \(i = 1, \ldots, q\)）。集合 E 是多项式凸的（引理 3），且公式 (3) 表明 \(\mathrm{pr}_{\Lambda'}(\mathrm{K}) \subset \mathrm{E}\)。

另一方面，设 \(c = (c_{\lambda})_{\lambda \in \Lambda'} \in \mathrm{E}\)；设 \(d = (d_{\lambda})_{\lambda \in \Lambda}\) 是 \(\mathbf{C}^{\Lambda}\) 中由 \(d_{\lambda} = c_{\lambda}\)（对 \(\lambda \in \Lambda'\)）与 \(d_{\lambda} = 0\)（对 \(\lambda \in \Lambda - \Lambda'\)）定义的元素。于是 (4) 蕴涵 \(d \in \mathrm{U}\)，从而 \(c \in \mathrm{pr}_{\Lambda'}(\mathrm{U})\)。这样，\(\mathrm{E} \subset \mathrm{pr}_{\Lambda'}(\mathrm{U})\)，证明完成。

引理 7. --- 设 \(n \geqslant 1\) 是一个整数，K 是 \(\mathbf{C}^{n}\) 的一个多项式凸紧子集。则 K 具有一个由多项式凸紧邻域组成的基本系。

存在一个紧多圆柱 (cf.~VAR, R1, p.~24) \(\Delta\)（\(\mathbf{C}^n\) 中的），它是 K 的一个邻域。由于 K 是多项式凸的，存在一个族 \((\mathrm{P}_i)_{i\in \mathrm{I}}\)（由 \(\mathbf{C}[\mathrm{X}_1,\dots ,\mathrm{X}_n]\) 的元素组成），以及一个由正实数组成的族 \((\mathrm{M}_i)_{i\in \mathrm{I}}\)，使得 K 是 \(z\in\) \(\Delta\) 中满足 \(|\mathrm{P}_i(z)|\leqslant \mathrm{M}_i\)（对一切 \(i\)）的 z 所成的集合（引理 3）。对 I 的每个有限子集 J 及每个 \(\varepsilon >0\)，用 \(\mathrm{K}_{\mathrm{J},\varepsilon}\) 表示 \(z\in \Delta\) 中满足 \(|\mathrm{P}_i(z)|\leqslant \mathrm{M}_i + \varepsilon\)（对 \(i\in \mathrm{J}\)）的 z 所成的集合。则每个集合 \(\mathrm{K}_{\mathrm{J},\varepsilon}\) 都是 K 的一个多项式凸紧邻域 (loc. cit.)，且诸集合 \(\mathrm{K}_{\mathrm{J},\varepsilon}\) 的交是 K。于是诸集合 \(\mathrm{K}_{\mathrm{J},\varepsilon}\) 构成 K 的多项式凸邻域的一个基本系 (TG, I, p.~60, th. 1)。

